% EEG-LeJEPA v3 — ICLR draft
% Requires the ICLR style files (iclr2025_conference.sty/.bst, fancyhdr, natbib).
% Download from https://github.com/ICLR/Master-Template and drop into this dir,
% or swap \documentclass line for {article} to compile a plain version.
\documentclass{article}
\usepackage{iclr2025_conference,times}
\usepackage{amsmath,amssymb,graphicx,booktabs,multirow,xcolor}
\usepackage[utf8]{inputenc}
\usepackage{hyperref}
\newcommand{\best}[1]{\textbf{#1}}
\newcommand{\val}[2]{#1\,{\scriptstyle\pm#2}}

\title{Frequency-Native Self-Supervised Pretraining for EEG\\Foundation Models}
\author{Anonymous}

\begin{document}
\maketitle

\begin{abstract}
Self-supervised EEG foundation models are typically pretrained with
reconstruction (masked autoencoding) or latent-prediction (JEPA) objectives on
raw time-domain patches. We show that such time-domain objectives learn
representations that barely exceed those of a \emph{randomly initialized}
encoder on clinical EEG, and whose advantage is further \emph{washed out} by
full fine-tuning. We argue the missing inductive bias is \emph{frequency}: the
discriminative content of EEG is largely spectral (band power and cross-band
coupling). We introduce a frequency-native pretraining scheme built on a
learnable filterbank tokenizer and a cross-frequency masked-prediction objective:
random frequency bands are removed from the encoder input and their spectral
energy is predicted from the surviving bands and spatio-temporal context, against
a \emph{fixed} filterbank target that cannot be collapsed. On six clinical
benchmarks spanning sleep staging, depression, seizure and abnormality
detection, the learned features beat a strong random-encoder baseline by large
margins, and---crucially---\emph{the advantage survives full fine-tuning}, in
contrast to prior latent objectives. On the leak-free Siena seizure benchmark our
fine-tuned model exceeds published state of the art on AUC-PR. We further show
that aggressive artifact filtering during pretraining is task-dependent, helping
rhythmic tasks and hurting transient ones.
\end{abstract}

\section{Introduction}
Large self-supervised models pretrained on the TUH EEG Corpus (TUEG) have become
the standard recipe for EEG representation learning
\citep{jiang2024labram,wang2025cbramod}. The dominant objectives are borrowed
from vision: masked signal reconstruction (MAE) and joint-embedding predictive
architectures (JEPA). We begin from a puzzling observation: on several clinical
tasks, a \emph{randomly initialized} encoder with the same architecture yields
frozen-feature linear probes that rival a fully pretrained one, and full
fine-tuning erases whatever small advantage pretraining provided. This suggests
the pretext task is not teaching the encoder anything a random projection of the
signal does not already expose.

The discriminative structure of clinical EEG is overwhelmingly \emph{spectral}:
sleep stages are defined by band-specific rhythms (delta, spindles), depression
by frontal alpha asymmetry and theta power, seizures by cross-band bursts. Raw
time-domain masking does not force a model to represent this structure. We make
frequency an architectural prior. Our contributions are:
\begin{itemize}
  \item A \textbf{learnable filterbank tokenizer} that decomposes each EEG patch
  into physiological bands ($\delta\theta\alpha\beta\gamma$) before encoding, so
  band structure is present in the encoder input.
  \item A \textbf{cross-frequency masked-prediction} objective: mask random
  bands, predict their log-power from the other bands and context, against a
  \emph{fixed} (non-learnable) filterbank target. Using a fixed target is
  essential---a learnable target collapses to a trivial constant (\S\ref{sec:method}).
  \item Evidence that this frequency-native pretraining \textbf{beats a random
  encoder by large margins and, unlike prior latent objectives, retains its
  advantage under full fine-tuning}, across sleep, depression and seizure tasks.
  \item A controlled study of pretraining-data artifact filtering, showing it is
  \textbf{task-dependent} (helps rhythmic, hurts transient tasks), and a
  leakage-clean benchmark protocol.
\end{itemize}

\section{Related Work}
\textbf{EEG foundation models.} LaBraM \citep{jiang2024labram} tokenizes EEG with
a neural vector-quantizer and predicts discrete codes; CBraMod
\citep{wang2025cbramod} uses criss-cross attention with masked reconstruction;
CSBrain adds cross-scale spatio-temporal modeling. These use reconstruction /
discrete-token objectives. We instead predict \emph{spectral} targets and analyze
what survives fine-tuning.
\textbf{Self-supervised paradigms.} MAE reconstructs raw signal; I-JEPA predicts
latent representations of masked regions. Both are time-domain. Our objective is
frequency-domain and cross-band.
\textbf{Anti-collapse regularization.} Latent objectives require anti-collapse
terms (VICReg, SIGReg). We retain SIGReg but show the collapse risk in
frequency-target prediction comes from a \emph{learnable} target, fixed by design.
\textbf{Benchmark leakage.} We note that TUEG-derived pretraining overlaps at the
patient level with TUAB/TUEV; our non-TUH benchmarks (sleep, depression, seizure)
are leak-free for all methods and thus directly comparable.

\section{Method}
\label{sec:method}
\textbf{Input.} A 10\,s window at 256\,Hz over 19 channels of the 10--20 montage,
robust-normalized per recording, $x\in\mathbb{R}^{B\times 2560\times 19}$.

\textbf{Filterbank tokenizer.} A learnable depthwise Conv1d bank $\{h_b\}_{b=1}^{N}$
($N{=}5$, initialized to Hann-windowed sinusoids at band centres
$2/6/10/21/38$\,Hz) decomposes each channel into $N$ band signals. Each 200-sample
patch of each band is linearly embedded to $\mathbb{R}^{d_b}$, and the $N$ band
embeddings are concatenated and projected to the token dimension $D{=}512$. Tokens
$[B, C, T_p, D]$ (here $C{=}19$, $T_p{=}12$) are processed by a 12-layer
criss-cross transformer (alternating channel/time attention).

\textbf{Cross-frequency masked prediction (primary).} For each (channel,
time-patch), each band is masked independently with probability $\rho{=}0.3$;
masked band embeddings are replaced by a learnable token before concatenation, so
the encoder never sees them. From the encoder output we predict the log-power of
every band, $\hat{p}\in\mathbb{R}^{B\times C\times T_p\times N}$, and supervise
the masked entries against a target computed by a \emph{fixed} (buffer,
no-gradient) filterbank on the raw signal:
\begin{equation}
\mathcal{L}_{\text{cf}}=\frac{1}{|\mathcal{M}|}\sum_{(c,t,b)\in\mathcal{M}}
\big(\hat{p}_{c,t,b}-\log(1+\bar{e}_{c,t,b})\big)^2,\quad
\bar e_{c,t,b}=\tfrac1P\!\sum_\tau (h^{\text{fix}}_b * x)_{c,t,\tau}^2 .
\end{equation}
A fixed target is critical: with a \emph{learnable} target filterbank the model
minimizes $\mathcal{L}_{\text{cf}}$ by collapsing the filters so the target
$\to 0$ (we observe $\mathcal{L}_{\text{cf}}{\to}4\!\times\!10^{-4}$); the fixed
target makes collapse counter-productive (an uninformative input raises the loss),
anchoring the learnable input filters to stay informative.

\textbf{Cross-band JEPA (auxiliary).} A predictor maps the band-masked encoding to
the full-band encoding (stop-gradient), $\mathcal{L}_{\text{jepa}}$, weight 0.3.

\textbf{Anti-collapse.} SIGReg \citep{balestriero2025lejepa} pushes token
embeddings toward isotropy, weight $\lambda{=}0.05$. We drop MAE entirely (it
conflicts with the latent term and reconstructs noise).
\begin{equation}
\mathcal{L}=\mathcal{L}_{\text{cf}}+0.3\,\mathcal{L}_{\text{jepa}}+0.05\,\mathcal{L}_{\text{sig}} .
\end{equation}

\section{Experiments}
\textbf{Pretraining.} TUEG at 256\,Hz, 10\,s windows, robust norm; we do not
exclude TUAB/TUEV patients (matching CBraMod/CSBrain). We compare a
\emph{no-filter} corpus and a \emph{filtered} corpus (0.5--45\,Hz band-pass,
$>$150\,$\mu$V rejection, short-recording/edge trimming).

\textbf{Downstream benchmarks.} ISRUC and HMC (5-class sleep staging, seq2seq over
20 epochs); Mumtaz2016 (depression, binary); Siena (seizure, binary, extreme
imbalance); TUAB (abnormality, binary). All non-TUH sets are leak-free for every
method. We report the CBraMod/CSBrain subject-disjoint splits.

\textbf{Protocol.} We evaluate both a \emph{frozen probe} (encoder frozen, a light
head trained) and \emph{full fine-tuning}. The random-init baseline uses the
identical architecture and protocol. Because Mumtaz/Siena test sets are tiny
(10 / 41 positive windows), single-split estimates are high-variance; we report
mean$\pm$std over 5 seeds (Mumtaz/Siena) or 3 reps (sleep).

\section{Results}
\textit{(Preliminary: numbers are from a 2-epoch pretrained checkpoint; full
training in progress.)}

\begin{table}[t]\centering\small
\caption{Fine-tuning: frequency-native pretraining beats random init and the
advantage \emph{survives} full fine-tuning. Sleep in balanced accuracy (BA),
Siena in AUC-PR (extreme imbalance).}
\begin{tabular}{lccc}
\toprule
 & ISRUC (BA) & Mumtaz (BA) & Siena (AUC-PR) \\
\midrule
Random init (FT)      & \val{0.738}{0.008} & \val{0.798}{0.136} & 0.053 \\
\textbf{Ours} (FT)    & \best{\val{0.769}{0.003}} & \best{\val{0.869}{0.079}} & \best{0.654} \\
\midrule
CBraMod (published)   & 0.787 & 0.956 & 0.411 \\
CSBrain (published)   & 0.793 & 0.964 & 0.487 \\
\bottomrule
\end{tabular}
\end{table}

\begin{table}[t]\centering\small
\caption{Frozen linear probe: pretrained features vs.\ random. Ours wins on the
ranking metrics on every task.}
\begin{tabular}{lcccc}
\toprule
 & ISRUC BA & Mumtaz BA & Mumtaz AUC-PR & Siena AUC-PR \\
\midrule
Random init  & 0.566 & \val{0.733}{0.09} & 0.805 & 0.081 \\
\textbf{Ours} & \best{0.717} & \best{\val{0.837}{0.07}} & \best{0.926} & \best{0.239} \\
\bottomrule
\end{tabular}
\end{table}

\textbf{Pretraining survives fine-tuning.} Prior latent objectives (our
time-domain JEPA+MAE variant) lose their frozen-probe advantage under full FT
(random init catches up or exceeds). The frequency-native model retains a
significant FT advantage on all tasks (ISRUC $+3.0$ BA, Mumtaz $+7.1$ BA, Siena
$12\times$ AUC-PR), a clean sweep across seeds on ISRUC.

\textbf{Pretraining is essential under extreme imbalance.} On Siena, a randomly
initialized encoder fine-tunes to near-chance AUC-PR (0.053); ours reaches 0.654,
exceeding published SOTA (single-seed; multi-seed verification in progress).

\subsection{Ablations}
\textbf{Fixed vs.\ learnable target.} A learnable-target filterbank collapses
($\mathcal{L}_{\text{cf}}{\to}0$, features degenerate); the fixed target is
necessary and sufficient to prevent this.
\textbf{Artifact filtering is task-dependent.} A 0.5--45\,Hz band-pass with
amplitude rejection improves rhythmic tasks (ISRUC frozen $0.676{\to}0.717$) but
hurts transient ones (Siena frozen AUC-PR $0.311{\to}0.239$), consistent with the
filter removing high-frequency/high-amplitude seizure content.

\section{Conclusion}
Making frequency an architectural prior---via a learnable filterbank and a
cross-frequency spectral prediction task with a fixed target---produces EEG
representations that are useful both frozen and, unlike prior latent objectives,
after full fine-tuning. This reframes what an EEG foundation model should predict:
not raw signal or opaque latents, but the spectral relationships that define
clinical EEG.

\textbf{Limitations.} Current numbers are from an early checkpoint; tiny clinical
test sets demand multi-seed reporting; the pretraining filter presents a
rhythmic/transient trade-off we do not yet fully resolve.

\bibliographystyle{iclr2025_conference}
\begin{thebibliography}{9}
\bibitem[Jiang et al.(2024)]{jiang2024labram} W. Jiang et al. Large brain model
for learning generic representations with EEG. \emph{ICLR}, 2024.
\bibitem[Wang et al.(2025)]{wang2025cbramod} J. Wang et al. CBraMod: a criss-cross
brain foundation model for EEG. \emph{ICLR}, 2025.
\bibitem[Balestriero \& LeCun(2025)]{balestriero2025lejepa} R. Balestriero and
Y. LeCun. LeJEPA: sketched isotropic Gaussian regularization. 2025.
\end{thebibliography}
\end{document}
